%!TEX TS-program = xelatex
%!TEX root = ../Topo2_WS15/topologie_2.tex
\RequirePackage{fix-cm} 
\documentclass[a4paper, twoside, headsepline, index=totoc,toc=listof, fontsize=10pt, cleardoublepage=empty, headinclude, DIV=12, BCOR=5mm, titlepage,draft]{scrartcl}
\usepackage{scrtime} % KOMA, Uhrzeit ermoeglicht

\usepackage{etoolbox}
\usepackage{letltxmacro}
\usepackage{ifthen}

%--Farbdefinitionen
%-- muss vor tikz geladen werden
\usepackage[usenames, table, x11names]{xcolor}
\definecolor{dark_gray}{gray}{0.45}
\definecolor{light_gray}{gray}{0.6}
\usepackage[final]{graphicx}

%--Zum Zeichnen
%-- muss vor polyglossia bzw. babel geladen werden
\usepackage{tikz}
\usepackage{tikz-cd}
\usetikzlibrary{external}
\tikzset{>=latex}
\usetikzlibrary{shapes,arrows.meta,intersections}
\usetikzlibrary{calc,3d}
\usetikzlibrary{decorations.pathreplacing,decorations.markings, decorations.pathmorphing}
\usetikzlibrary{angles}

%-- Konfiguration von tikzexternalize
\tikzexternalize[prefix=tikz/,up to date check=diff]
\pgfkeys{/pgf/images/include external/.code=\includegraphics{#1}}
\tikzset{external/system call={xelatex \tikzexternalcheckshellescape -halt-on-error -interaction=batchmode --shell-escape -jobname "\image" "\texsource"}}


%-- tikzexternalize fuer tikzcd deaktivieren, da inkompatibel
\AtBeginEnvironment{tikzcd}{\tikzexternaldisable}
\AtEndEnvironment{tikzcd}{\tikzexternalenable}

%-- um Inkompatibilitaeten von quotes und polyglossia bzw. babel zu vermeiden
\tikzset{
  every picture/.append style={
    execute at begin picture={\shorthandoff{"}},
    execute at end picture={\shorthandon{"}}
  }
}
\usetikzlibrary{quotes}
\usepackage{pgfplots}
\usepgfplotslibrary{colormaps}
% \pgfplotsset{compat=1.12}



%-- Mathesymbole etc.
\usepackage{mathtools} % beinhaltet amsmath
\mathtoolsset{showonlyrefs, centercolon,showmanualtags}
\newtagform{brackets}[\textbf]{[}{]}
\usetagform{brackets}
\usepackage{fix-cm}
% \usepackage{wasysym} % zusätzliche Symbole
\usepackage{amssymb} % zusätzliche Symbole
% \usepackage{stmaryrd} % für Blitz
\usepackage{nicefrac} % schräge Brüche
\usepackage{faktor}
\newcommand{\Faktor}[1]{\faktor[\textstyle]{#1}}
\usepackage{xfrac}
\usepackage{cancel}
\usepackage{mathdots} % Verbesserung von Punkten wie zB \ldots
\usepackage[bb=px]{mathalfa}

%-- einzelne Symbole aus dem mathx-package
% \DeclareFontFamily{U}{mathx}{\hyphenchar\font45}
% \DeclareFontShape{U}{mathx}{m}{n}{<-> mathx10}{}
% \DeclareSymbolFont{mathx}{U}{mathx}{m}{n}
% \DeclareMathSymbol{\bigplus}{1}{mathx}{"90}

%-- 
\usepackage{silence}
\WarningFilter{latexfont}{Size substitutions with differences}
\WarningFilter{latexfont}{Font shape `U/bbold/m/n' in size}

\DeclareSymbolFont{bbold}{U}{bbold}{m}{n}
\DeclareSymbolFontAlphabet{\mathbbold}{bbold}
\newcommand{\ind}{\mathbbold{1}} % charakteristische-Funktion-Eins

\def\mathul#1#2{\color{#1}\underline{{\color{black}#2}}\color{black}} %farbiges Untersteichen im Mathe-Modus
\renewcommand{\le}{\leqslant}
\renewcommand{\ge}{\geqslant}

%-- Underbrace als Befehl in LaTeX-Syntax (und ohne Spacingprobleme mit nachfolgenden Operatoren...)
\newcommand{\Underbrace}[2]{{\underbrace{#1}_{#2}}}
\newcommand{\Underbracket}[2]{{\underbracket[0.7pt][2pt]{#1}_{#2}}}

%-- Alles was mit Schrift und XeTeX zu tun hat
\usepackage{palatino}
\usepackage[euler-digits]{eulervm}
\usepackage[no-math]{fontspec}
\usepackage{polyglossia} % moderner babel-ersatz
\setmainlanguage[spelling=new,babelshorthands=true]{german}
\shorthandoff{"}
\setotherlanguage{english}
\defaultfontfeatures{Mapping=tex-text, WordSpace={1.2}, Ligatures={Required,Common,Contextual}} %

%-- Alle hier aufgeführten Schriftarten sind in jeder TeXlive-Installation vorhanden!!
%-- Hauptschriftart festlegen
% \setmainfont{texgyrepagella}[Extension=.otf,UprightFont=*-regular,BoldFont=*-bold,ItalicFont=*-italic,BoldItalicFont=*-bolditalic,ItalicFeatures={Style=Historic}]
\setmainfont{TeXGyrePagellaX}[Extension=.otf,UprightFont=*-Regular,BoldFont=*-Bold,ItalicFont=*-Italic,BoldItalicFont=*-BoldItalic,ItalicFeatures={Style=Historic},Ligatures={Required,Common,Contextual,Historic}]
% \setmainfont[BoldFont={* Bold}, ItalicFont={* Light Italic},SmallCapsFont={Linux Libertine O}, SmallCapsFeatures={Letters=SmallCaps}]{Source Sans Pro}


%-- Serifenlose Schriftart festlegen
% \setsansfont{Roboto}[Scale=MatchUppercase,Extension=.ttf,UprightFont=*-Regular,BoldFont=*-Bold,ItalicFont=*-RegularItalic,BoldItalicFont=*-MediumItalic]
% \setsansfont{RobotoSlab}[Scale=MatchUppercase,Extension=.ttf,UprightFont=*-Regular,BoldFont=*-Bold,ItalicFont=*-Thin,SmallCapsFont=*-Regular]
% \setsansfont{FiraSans}[Scale=MatchUppercase,Extension=.otf, UprightFont=*-Regular, BoldFont=*-SemiBold, ItalicFont=*-RegularItalic, BoldItalicFont=*-MediumItalic]
% \setsansfont{SourceSansPro}[Scale=MatchUppercase,Extension=.otf,UprightFont=*-Regular,BoldFont=*-Semibold,ItalicFont=*-LightIt,BoldItalicFont=*-SemiboldIt]
% \setsansfont{texgyreadventor}[Scale=MatchUppercase, Extension=.otf, UprightFont=*-regular, BoldFont=*-bold, ItalicFont=*-italic,BoldItalicFont=*-bolditalic]
% \setsansfont{Kurier}[Scale=MatchUppercase, Extension=.otf, UprightFont=*Medium-Regular, BoldFont=*Heavy-Regular, ItalicFont=*Medium-Italic,BoldItalicFont=*Heavy-Italic]
% \setsansfont{Iwona}[Scale=MatchUppercase, Extension=.otf, UprightFont=*Medium-Regular, BoldFont=*Heavy-Regular, ItalicFont=*Medium-Italic,BoldItalicFont=*Heavy-Italic]
% \setsansfont{LinBiolinum}[Scale=MatchUppercase, Extension=.otf, UprightFont=*_R, BoldFont=*_RB, ItalicFont=*_RI,BoldItalicFont=*_RBO]
% \setsansfont{texgyrebonum}[Scale=MatchUppercase,Extension=.otf, UprightFont=*-regular, BoldFont=*-bold, ItalicFont=*-italic, BoldItalicFont=*-bolditalic]
\setsansfont{texgyreadventor}[Scale=MatchUppercase,Extension=.otf, UprightFont=*-regular, BoldFont=*-bold, ItalicFont=*-italic, BoldItalicFont=*-bolditalic]



% \newfontfamily{\geometric}{texgyreheros}[Scale=MatchLowercase, Extension=.otf, UprightFont=*-regular, BoldFont=*-bold, ItalicFont=*-italic]
% \addtokomafont{disposition}{\geometric}
\setmonofont{SourceCodePro}[Scale=0.9,Extension=.otf,UprightFont=*-Regular, BoldFont=*-Semibold, ItalicFont=*-Light]
\usepackage{xltxtra}
\usepackage{fontawesome}
\usepackage[final]{microtype}


\usepackage[draft=false]{scrlayer-scrpage} % wie fancyhdr, nur optimiert auf KOMA-Skript, leicht andere Syntax

\flushbottom


%--Mixed
\usepackage[shortlabels,inline]{enumitem}
\setlist[itemize,1]{label=\faCaretRight}
\setlist[enumerate]{font=\bfseries}
\setlist[description]{font=\normalfont\bfseries}
\usepackage[german=quotes]{csquotes}
\usepackage{makeidx}
\usepackage{wrapfig}
\usepackage{float}
\usepackage[margin=10pt, font=small, labelfont={sf, bf}, format=plain, indention=1em]{caption}
\captionsetup[wrapfigure]{name=Abb. }
\usepackage{stackrel}


%--Unterstreichung
\usepackage[normalem]{ulem}
\setlength{\ULdepth}{1.8pt}

%--Indexverarbeitung
\newcommand{\bet}[1]{\emph{#1}}
\newcommand{\Index}[1]{\bet{#1}\index{#1}}
\makeindex
\setindexpreamble{{\noindent\sffamily\small Die \emph{Seitenzahlen} sind mit Hyperlinks versehen und somit anklickbar}  \par \bigskip}
\renewcommand{\indexpagestyle}{scrheadings}


%--Marginnote & todonotes
\deffootnote[1.5em]{1.5em}{1.5em}{\textsuperscript{\thefootnotemark}\ }
\usepackage[fulladjust]{marginnote}
\renewcommand*{\marginfont}{\color{dark_gray} \itshape \footnotesize}
\usepackage[textsize=small]{todonotes}
\usepackage{ragged2e}
\renewcommand*{\raggedleftmarginnote}{\RaggedLeft}
\renewcommand*{\raggedrightmarginnote}{\RaggedRight}

%--Todonotes mit tikz/externalize kompatibel machen
\LetLtxMacro{\oldtodo}{\todo}
\renewcommand{\todo}[2][]{\tikzexternaldisable\oldtodo[#1]{#2}\tikzexternalenable}
\LetLtxMacro{\oldmissingfigure}{\missingfigure}
\renewcommand{\missingfigure}[2][]{\tikzexternaldisable\oldmissingfigure[{#1}]{#2}\tikzexternalenable}

%--Konfiguration von Hyperref pdfstartview=FitH, 
\usepackage[hidelinks, pdfpagelabels,  bookmarksopen=true, bookmarksnumbered=true, linkcolor=black, urlcolor=SkyBlue2, plainpages=false,pagebackref, citecolor=black, hypertexnames=true, pdfauthor={Jannes Bantje}, pdfborderstyle={/S/U}, linkbordercolor=SkyBlue2, colorlinks=false,final]{hyperref}


\newcommand{\appendLink}[1]{#1\,\faExternalLink}
\newcommand{\hrefsym}[2]{\href{#1}{\texttt{\appendLink{#2}}}}
\newcommand{\hrefsymmail}[2]{\href{#1}{\texttt{\faEnvelopeO\,#2}}}
\renewcommand{\url}[1]{\hrefsym{#1}{\nolinkurl{#1}}}



% -- QR-Codes
\usepackage{qrcode} %-- hinter hyperref laden!

%--Römische Zahlen
\newcommand{\RM}[1]{\MakeUppercase{\romannumeral #1{}}}

%%--Abkürzungen etc.
\newcommand{\light}{\text{\Large $\lightning$}}

%-- Definitionen von weiteren Mathe-Befehlen

\DeclareMathOperator{\im}{im} %Bild
\DeclareMathOperator{\id}{id} %identische Abbildung
\DeclareMathOperator{\conj}{conj} %Konjugation
\DeclareMathOperator{\sgn}{sgn} %Signum
\DeclareMathOperator{\End}{End} %Endomorphismen
\DeclareMathOperator{\Hom}{Hom} %Homomorphismen
\DeclareMathOperator{\Iso}{Iso} %Isomorphismen
\DeclareMathOperator{\Span}{span} %Span

\DeclareMathOperator{\grad}{grad} %Gradient
\DeclareMathOperator{\dive}{div} %Gradient
\DeclareMathOperator{\rot}{rot} %Rotation


\DeclareMathOperator{\Char}{char} %Charakteristik
\DeclareMathOperator{\Aut}{Aut} %Automorphismen

% -- Zum Finetuning von Befehlen
\makeatletter
\newcommand{\raisemath}[1]{\mathpalette{\raisem@th{#1}}}
\newcommand{\raisem@th}[3]{\raisebox{#1}{$#2#3$}}
\makeatother

\newcommand{\sing}{{\raisemath{1.1pt}{\scriptscriptstyle\mathrm{sing}}}}
\newcommand{\pt}{\mathrm{pt}}
\newcommand{\op}{\mathrm{op}}
\DeclareMathOperator{\Sp}{Sp}
\DeclareMathOperator{\Keg}{Keg}
\newcommand{\slashedi}{i\hspace{-3.5pt}/}
\newcommand{\cupp}{\smallsmile}
\newcommand{\capp}{\smallfrown}
\DeclareMathOperator*{\colim}{colim}


\newcommand{\D}{\ensuremath{\mathrm{D}\mkern-1.5mu}}
\newcommand{\mathd}{\ensuremath{\mathrm{d}\mkern-1.0mu}}
\newcommand{\Tmap}{\ensuremath{\mathrm{T}\mkern-0.85mu}}
\newcommand{\opL}{\ensuremath{\mathrm{L}\mkern-0.6mu}}


\DeclareMathOperator{\supp}{supp} %Träger


% -- Befehle für FunkAna und Operatoralgebren
\newcommand{\tr}{\mathrm{tr}}
\newcommand{\w}{\mathrm{w}}
\newcommand{\sa}{\mathrm{sa}}
\newcommand{\so}{\mathrm{s.o.}}
\newcommand{\weakT}[1]{\ensuremath{\mathcal{T}_{#1}^{\mkern+1.0mu\text{\raisebox{0.4ex}{$\mathrm{w}$}}}}}
\newcommand{\weakTstar}[1]{\ensuremath{\mathcal{T}_{#1}^{\mkern+1.0mu\text{\raisebox{0.4ex}{$\mathrm{w}$}}^*}}}
\newcommand{\finSub}{\subset\mkern-0.7mu \subset}
\DeclareMathOperator{\Inv}{Inv}
\DeclareMathOperator{\dist}{dist}
\newcommand{\simm}{{\hspace{-1.6pt}\raisemath{0.5pt}{\sim}}}
\DeclareMathOperator{\ev}{ev}
\DeclareMathOperator{\Alg}{Alg}


% -- Kategorien
\newcommand{\TOP}{\textsc{Top}}
\newcommand{\HTOP}{\textsc{HTop}}
\newcommand{\VR}{\textsc{VR}}
\newcommand{\MOD}{\textsc{Mod}}
\newcommand{\MONOIDE}{\textsc{Monoide}}
\newcommand{\MENGEN}{\textsc{Mengen}}
\newcommand{\MAN}{\textsc{Man}}
\newcommand{\GRUPPEN}{\textsc{Gruppen}}
\newcommand{\ABELGRUPPEN}{\textsc{Abel.Gruppen}}
\newcommand{\KAT}{\textsc{Kat}}
\newcommand{\FUN}{\textsc{Fun}}
\newcommand{\SIMP}{\textsc{Simp}}
\newcommand{\VEKT}{\textsc{Vekt}}

\DeclareMathOperator{\cell}{cell}
\DeclarePairedDelimiter{\homologieklasse}{\llbracket}{\rrbracket}
\newcommand{\rand}[1]{\ensuremath{\partial^{\scriptscriptstyle #1}}}


% -- theorem packages
\usepackage{amsthm}
\usepackage{thmtools,thm-restate}
\renewcommand{\listtheoremname}{Übersicht aller Aussagen}

% -- Theoreme als PDF-Lesezeichen
\usepackage{bookmark}
\bookmarksetup{open,numbered}
\makeatletter
\newcommand*{\theorembookmark}{%
  \bookmark[
    dest=\@currentHref,
    rellevel=1,
    keeplevel,
  ]{%
    \thmt@thmname\space\csname the\thmt@envname\endcsname
    \ifx\thmt@shortoptarg\@empty
    \else
      \space(\thmt@shortoptarg)%
    \fi
  }%
}   
\makeatother

% -- Definition der einzelnen Umgebungen
\declaretheoremstyle[%
	headfont=\sffamily\bfseries,
	notefont=\normalfont\sffamily\scshape,
	bodyfont=\normalfont,
	headformat=\NUMBER\ \NAME\NOTE,
	headpunct=.,
	postheadspace=1em,
	spaceabove=15pt,spacebelow=10pt,
	postheadhook=\theorembookmark]%
{mainstyle}
\declaretheoremstyle[%
	headfont=\sffamily\bfseries,
	notefont=\normalfont\sffamily\scshape,
	bodyfont=\normalfont,
	headformat=\NUMBER\NAME\NOTE,
	headpunct=.,
	postheadspace=1em,
	spaceabove=15pt,spacebelow=10pt,
	postheadhook=\theorembookmark]%
{mainstyle_unnumbered}
\declaretheoremstyle[%
	headfont=\sffamily\bfseries,
	notefont=\normalfont\sffamily\scshape,
	bodyfont=\normalfont,
	headformat=swapnumber,
	headpunct=.,
	postheadspace=1em,
	spaceabove=15pt,spacebelow=10pt,
	postheadhook=\theorembookmark,
	qed=\qedsymbol]%
{mainstyleB}
\declaretheorem[name=Definition,parent=section,style=mainstyle]{definition}
\declaretheorem[name=Definition,numbered=no,style=mainstyle_unnumbered]{definition*}
\declaretheorem[name=Theorem,sharenumber=definition,style=mainstyle]{theorem}
\declaretheorem[name=Theorem,numbered=no,style=mainstyle_unnumbered]{theorem*}
\declaretheorem[name=Proposition,sharenumber=definition,style=mainstyle]{proposition}
\declaretheorem[name=Lemma,sharenumber=definition,style=mainstyle]{lemma}
\declaretheorem[name=Satz,sharenumber=definition,style=mainstyle]{satz}
\declaretheorem[name=Korollar,sharenumber=definition,style=mainstyle]{korollar}
\declaretheorem[name=Korollar,sharenumber=definition,style=mainstyleB]{korollarB}
\declaretheorem[name=Frage,numbered=no,style=mainstyle_unnumbered]{frage}
\declaretheorem[name=Erinnerung,sharenumber=definition,style=mainstyle]{erinnerungA}


% -- Beweise
\declaretheoremstyle[headfont=\bfseries\scshape,bodyfont=\normalfont,headpunct=:,postheadspace=1em,spacebelow=12pt,spaceabove=2pt,qed=\qedsymbol]{beweise}
\declaretheoremstyle[headfont=\sffamily\bfseries,bodyfont=\normalfont,headpunct=:,postheadspace=1em,spacebelow=10pt,spaceabove=10pt]{bemerkungen}
\declaretheorem[name=Beweis,numbered=no,style=beweise]{beweis}
\declaretheorem[name=Bemerkung,sharenumber=definition,style=mainstyle]{bemerkung}
\declaretheorem[name=Beispiel,sharenumber=definition,style=mainstyle]{beispiel}
\declaretheorem[name=Übung,numbered=no,style=bemerkungen]{uebung}
\declaretheorem[name=Erinnerung,numbered=no,style=bemerkungen]{erinnerung}


%--Skalarprodukt
\DeclarePairedDelimiterX\sprod[2]{\langle}{\rangle}{#1\,\delimsize\vert\,#2}
\DeclarePairedDelimiterX\skal[2]{\langle}{\rangle}{#1\,,\,#2}

%--Betrag, Gaußklammer
\DeclarePairedDelimiter{\abs}{\lvert}{\rvert}
\DeclarePairedDelimiter{\floor}{\lfloor}{\rfloor}
\DeclarePairedDelimiter{\ceil}{\lceil}{\rceil}

%--Norm
\DeclarePairedDelimiter\norm{\Vert}{\Vert}


%--Umklammern
\DeclarePairedDelimiter\enbrace{(}{)}
\DeclarePairedDelimiter\benbrace{[}{]}
\DeclarePairedDelimiter\lenbrace{<}{>}
\newcommand{\ssbrace}[1]{{\scriptscriptstyle\enbrace{#1}}}

%--Mengen
\newcommand\SetSymbol[1][]{\nonscript\:#1\vert\allowbreak\nonscript\:\mathopen{}}
\providecommand\given{} % to make it exist
\DeclarePairedDelimiterX\set[1]\{\}{\renewcommand\given{\SetSymbol[\delimsize]}#1}

%--Abbildungsdefinition
\newcommand{\mapdef}[5]{%
	\[
		\begin{array}{rcl}
			\textstyle #1 &\xrightarrow{\minwidthbox{#5}{2em}} & \textstyle #2 \\[0.5ex]
			\textstyle #3 &\xmapsto{\minwidthbox{\mbox{ }}{2em}} & \textstyle #4
		\end{array}
	\]
}

%--modifiziertes Stackrel
\newcommand{\StackText}[2]{\stackrel{\mbox{\scriptsize #1}}{#2}}

%--Befehl für eine Box mit Mindestweite
\DeclareRobustCommand{\minwidthbox}[2]{%
  \ifmmode
    \expandafter\mathmakebox
  \else
    \expandafter\makebox
  \fi
  [\ifdim#2<\width\width\else#2\fi]{#1}%
}

%--Befehl zum Einkreisen
\newcommand*\circled[1]{\tikzexternaldisable\tikz[baseline=(char.base)]{\node[shape=circle,draw,inner sep=2pt] (char) {#1};}\tikzexternalenable}


%--Differential
\newcommand{\diff}[2]{\ensuremath{\frac{{\partial #1}}{{\partial #2}} }}
\newcommand{\diffd}[2]{\ensuremath{\frac{\mathd #1}{\mathd #2} }}

%--Inhaltsverzeichnis
\usepackage[tocindentauto]{tocstyle}
\usetocstyle{KOMAlike}	
\shorthandon{"}
		
%--Punkte (als letztes laden)
\usepackage{ellipsis}

\newcommand{\fach}{Operatoralgebren}
\newcommand{\semester}{WiSe 2015}
\newcommand{\homepage}{https://wwwmath.uni-muenster.de/u/wilhelm.winter/wwinter/operatoralgebren.html}

\newcommand{\prof}{Prof.\ Dr.\ Wilhelm Winter}
% \publishers{\scalebox{14}{$K$}}
\providecommand{\verfasser}{Jannes Bantje}

%--Konfiguration von scrheadings
\setheadsepline{1pt}[\color{light_gray}]
\pagestyle{scrheadings}
%\fancyhf[H, F]{}
\clearscrheadfoot

\providecommand{\shortFach}{\fach}
\lehead{ \includegraphics[height=0.5 cm,keepaspectratio]{../!config/Bilder/Logo_WWU_Muenster_light_gray.pdf}}
\rehead{\rule{0cm}{0.5cm}\footnotesize \sffamily \color{light_gray} \verfasser -- Mitschrift \shortFach}
\lohead{\rule{0cm}{0.5cm} \footnotesize \sffamily \color{light_gray} Stand: \today \; \thistime[:]}
\rohead{\includegraphics[height=0.5 cm,keepaspectratio]{../!config/Bilder/fb10logo_gray.pdf}}


\ofoot[{ \color{dark_gray} \LARGE \sffamily \thepage}]{{ \color{dark_gray} \LARGE \sffamily \thepage}} %hier wir auch der plain Stil bearbeitet!
\automark{section}
\ifoot{ \color{dark_gray} \small \leftmark}

%--Metadaten
\author{\verfasser}
\titlehead{\includegraphics[height=1.5cm, keepaspectratio]{../!config/Bilder/Logo_WWU_Muenster.pdf}%
\hfill \includegraphics[height=1.3cm, keepaspectratio]{../!config/Bilder/fb10logo.pdf}}
\title{Skript \fach}
\subtitle{Mitschrift der Vorlesung  \enquote{\fach} von \prof}

\publishers{
Aktuelle Version verfügbar bei: \bigskip\\
\normalsize
\begin{minipage}[t]{0.45\textwidth}
	\centering{
	\qrcode[height=3.5cm, version=4]{https://github.com/JaMeZ-B/latex-wwu}
	\medskip\\
	\includegraphics[height=0.4cm, keepaspectratio]{../!config/Bilder/github_octo.pdf} 
	\includegraphics[height=0.4cm, keepaspectratio]{../!config/Bilder/GitHub_Logo.pdf} (inklusive Sourcecode)\\
	{\small\url{https://github.com/JaMeZ-B/latex-wwu}}}
\end{minipage}
\hfill
\begin{minipage}[t]{0.45\textwidth}
	\centering{
	\qrcode[height=3.5cm, version=4]{btsync://B6WH2DISQ5QVYIRYIEZSF4ZR2IDVKPN3I?n=latex_share}
	\medskip\\
	\includegraphics[height=0.4cm, keepaspectratio]{../!config/Bilder/bt_sync_logo.pdf} 
	{\large \textbf{Bittorrent} Sync} \\
	\texttt{B6WH2DISQ5QVYIRYIEZSF4ZR2IDVKPN3I}
	}
\end{minipage}
}


\begin{document}
\pagenumbering{Roman}
\maketitle
\begin{abstract}
\section*{Vorwort --- Mitarbeit am Skript}
Dieses Dokument ist eine Mitschrift aus der Vorlesung \enquote{\fach, \semester}, gelesen von \prof. Der Inhalt entspricht weitestgehend dem Tafelanschrieb. Für die
Korrektheit des Inhalts übernehme ich keinerlei Garantie! Für Bemerkungen und Korrekturen -- und seien es nur Rechtschreibfehler -- bin ich sehr dankbar. 
Korrekturen lassen sich prinzipiell auf drei Wegen einreichen: 
\begin{itemize}
	\item Persönliches Ansprechen in der Uni, Mails an \hrefsymmail{mailto:\mail}{\mail} (gerne auch mit annotieren PDFs) 
	\item \emph{Direktes} Mitarbeiten am Skript: Den Quellcode poste ich auf GitHub (siehe oben), also stehen vielfältige Möglichkeiten der Zusammenarbeit zur Verfügung:
	Zum Beispiel durch Kommentare am Code über die Website und die Kombination Fork + Pull Request. Wer sich verdient macht oder ein Skript zu einer Vorlesung, die 
	ich nicht besuche, beisteuern will, dem gewähre ich gerne auch Schreibzugriff.
	
	Beachten sollte man dabei, dass dazu ein Account bei \url{github.com} notwendig ist, der allerdings ohne Angabe von persönlichen Daten angelegt werden kann. 
	Wer bei GitHub (bzw. dem zugrunde liegenden Open-Source-Programm \enquote{\texttt{git}}) -- verständlicherweise -- Hilfe beim Einstieg braucht, dem helfe ich gerne 
	weiter. Es gibt aber auch zahlreiche empfehlenswerte Tutorials im Internet.\footnote{zB. \url{https://try.github.io/levels/1/challenges/1}, ist auf Englisch, aber dafür 
	interaktives LearningByDoing}
	\item \emph{Indirektes} Mitarbeiten: \TeX-Dateien per Mail verschicken. 
	
	Dies ist nur dann sinnvoll, wenn man einen ganzen Abschnitt ändern möchte (zB. einen alternativen Beweis geben), da ich die Änderungen dann per Hand einbauen muss! Ich freue mich aber auch über solche Beiträge!
\end{itemize}
\section*{Vorlesungshomepage}
{\centering 
\begin{minipage}[c][][c]{\textwidth}
	\centering \qrcode[height=3cm]{\homepage} \medskip\\
	\footnotesize \url{\homepage}
\end{minipage}
\par}

\section*{Anmerkung}
Innerhalb dieser Mitschrift wird man öfter den Ausdruck \enquote{Warum?} finden. Dies sind vom Dozenten bewusst weggelassene Details, die zu verstärktem Mitdenken beim Lesen 
animieren sollen. Oftmals sind dies schon aus vorherigen Semestern bekannte Sachverhalte. Nur an wenigen Stellen habe ich die fehlenden Details hinzugefügt.
\end{abstract}

\tableofcontents
\cleardoubleoddemptypage

\pagenumbering{arabic}
\setcounter{page}{1}
\setcounter{footnote}{0}

\section{Spektraltheorie in Banachalgebren} % (fold)
\label{sec:1}

\begin{definition}[{name=[Banachalgebra]}]
	Eine \Index{Banachalgebra} ist eine vollständige normierte Algebra. Für die Multiplikation gilt also $\norm{a b} \le \norm{a} \cdot \norm{b}$
	Wir betrachten in dieser Vorlesung fast ausschließlich Algebren über $\mathbb{C}$. Eine Banachalgebra heißt \Index{unital}, falls sie ein Einselement $\ind$ mit 
	$\norm{\ind} =1$ enthält. 
\end{definition}

\begin{definition}[{name=[Menge der invertierbaren Elemente und Spektrum]}]
	Für eine unitale $\mathbb{C}$-Algebra $A$ setzen wir 
	\begin{align}
		\Inv(A) &:= \set{a \in A  \text{ invertierbar} } \\
		\sigma(a) &:=\sigma_A(a) := \set[\big]{\lambda \in \mathbb{C} \given \lambda \cdot \ind - a \notin \Inv(A)} \marginnote{das Vorzeichen der Differenz spielt für die Definition der Menge keine Rolle}
	\end{align}
	$\sigma(a)$ heißt \Index{Spektrum} von $a$. 
\end{definition}

\begin{bemerkung}[{name=[Invertierbarkeit von $1 -ab$]},label=bem:13]
	Sei $A$ wie oben und $a,b \in A$. Dann gilt
	\[
		\ind -ab \in \Inv(A) \iff \ind -ba \in \Inv(A)
	\]
\end{bemerkung}
\begin{beweis}
	Durch simples Nachrechnen verifiziert man, dass $(\ind-ba)^{-1} = \ind+ b(\ind- ab)^{-1} a$ gilt.
\end{beweis}

\begin{satz}[label=satz14_poly,{name=[Verträglichkeit des Spektrums mit Polynomen]}]
	Sei $A$ eine unitale $\mathbb{C}$-Algebra und $a \in A$ mit $\sigma(a)\not= \emptyset$. Sei weiter $p \in \mathbb{C}[Z]$ ein Polynom. Dann gilt 
	\[
		\sigma\enbrace[\big]{p(a) }= p \enbrace[\big]{\sigma(a)}
	\]
\end{satz}
\begin{beweis}
	Die Aussage ist trivial, falls $p$ konstant ist. Sei also $p$ nicht konstant. Zu $\mu \in \mathbb{C}$ existieren $\lambda_0, \ldots , \lambda_n \in \mathbb{C}$ mit 
	$\lambda_0 \not= 0$ und $p(z) -\mu = \lambda_0 \cdot (z- \lambda_1) \cdot  \ldots \cdot (z- \lambda_n)$, also
	\[
		p(a) - \mu = \lambda_0 \cdot (a-\lambda_1) \cdot  \ldots \cdot (a- \lambda_n)
	\]
	Damit gilt: $p(a) -\mu$ ist genau dann invertierbar, wenn $a -\lambda_1, \ldots , a-\lambda_n $ invertierbar sind. Also 
	\begin{align}
		\mu \in \sigma \enbrace[\big]{p(a)} &\iff \lambda_i \in \sigma(a) \text{ ist für ein } i \in \set{1,\ldots ,n} \\
		&\iff p(\lambda) -\mu =0 \text{ für ein } \lambda \in \sigma(a) \\
		&\iff p(\lambda)=\mu \text{ für ein } \lambda \in \sigma(a)   
	\end{align}
	Somit folgt dann $\sigma \enbrace*{p(a)} = p \enbrace*{\sigma(a)}$.
\end{beweis}

\begin{satz}[label=satz15_neumann,{name=[Neumannsche Reihe]}]
	Sei $A$ eine unitale Banachalgebra und $a \in A$. Falls $\norm{a}<1$, so ist $\ind-a$ invertierbar und 
	\[
		(\ind-a)^{-1} = \sum_{n=0}^{\infty} a^n \tag*{(\Index{Neumannsche Reihe})}
	\]
\end{satz}
\begin{beweis}
	Die Folge der Partialsummen $\enbrace[\big]{\sum_{n=0}^{k} a^n}_{k \in \mathbb{N}} \subset A$ ist eine Cauchyfolge (warum?), sodass wegen Vollständigkeit die unendliche 
	Summe existiert. Es gilt dann
	\begin{align}
		(\ind-a) \enbrace*{\sum_{n=0}^{\infty} a^n} \xleftarrow{k \to \infty}(\ind-a) \enbrace*{\sum_{n=0}^{k} a^n} = \ind-a^{k+1} \xrightarrow{k \to \infty} \ind
	\end{align}
	Die Summe ist also ein Rechtsinverses. Genauso folgt, dass sie auch ein Linksinverses ist.
\end{beweis}

\begin{korollar}[{name=[Invertierbare Elemente bilden offene Menge]},label=inv_offen]
	Für jede unitale Banachalgebra $A$ ist $\Inv(A) \subset A$ offen.
\end{korollar}
\begin{beweis}
	Falls $a \in \Inv(A)$ und $\norm{b-a} < \norm{a^{-1}}^{-1}$, so gilt $\norm{b a^{-1} -\ind} \le \norm{b-a} \cdot \norm{a^{-1}}<1$.  Also ist auch
	$-\enbrace*{\ind- \enbrace*{ba^{-1} -\ind}} = ba^{-1}$ invertierbar nach \autoref{satz15_neumann} und damit auch $b \in \Inv(A)$, denn $b$ ist ein Linksinverses für 
	$a^{-1}(ba^{-1})^{-1} \in \Inv(A)$.
\end{beweis}

\begin{definition}[{name=[Spektralradius]}]
	Sei $A$ eine unitale Banachalgebra. Für $a \in A$ definieren wir den \Index{Spektralradius} durch
	\[
		r(a) := \sup_{\lambda \in \sigma(a)} \abs{\lambda} 
	\] 
\end{definition}

\begin{proposition}[{name=[Spektrum ist kompakt und Spektralradius durch Norm beschränkt]},label=spec_kompakt]
	Seien $a \in A$ wie oben. Dann ist das Spektrum $\sigma(a)$ kompakt und $r(a) \le \norm{a}$.
\end{proposition}
\begin{beweis}
	Angenommen $\abs{\lambda} > \norm{a}$, dann gilt $\norm{\lambda^{-1} a} < 1$, also ist $\ind- \lambda^{-1} a$ invertierbar nach \autoref{satz15_neumann} und 
	damit auch $\lambda -a \iff \lambda \notin \sigma(a)$. Insgesamt folgt, dass $r(a) \le \norm{a}$ ist. 	
	Es ist $\mathbb{C} \setminus \sigma(a) \subset \mathbb{C}$ offen, da $\Inv(A) \subset A$ offen ist. Damit muss $\sigma(a)$ nach Heine-Borel kompakt sein.
\end{beweis}

\begin{satz}[{name=[In unitalen Banachalgebren sind die Spektren nie leer]},label=spec_nichtleer]
	Sei $A$ eine unitale Banachalgebra. Für jedes $a \in A$ ist $\sigma(a)$ nichtleer. 
\end{satz}
\begin{beweis}
	Falls $\abs{\lambda}> \norm{a}$, so existiert $\enbrace*{\ind- \lambda^{-1} a}^{-1}$ und $(\lambda-a)^{-1}$ mit der gleichen Begründung wie eben. Es gilt dann
	\begin{align}
		(\lambda-a)^{-1} = \lambda^{-1} \enbrace*{1- \lambda^{-1} a}^{-1} \stackrel{\mbox{\scriptsize\ref{satz15_neumann}}}{=} \lambda^{-1} \enbrace*{\sum_{n=0}^{\infty} 
		\lambda^{-n}a^n}
	\end{align}
	und weiter
	\[
		\norm{(\lambda-a)^{-1}} \le \abs*{\lambda^{-1}} \cdot \sum_{n=0}^{\infty} \enbrace*{\abs*{\lambda^{-1}} \cdot \norm{a}}^n  \xrightarrow{\abs{\lambda} \to \infty} 0 
		\label{19:eq:1} \tag{$*$}
	\]
	da man die Summe als geometrische Reihe mit $2$ abschätzen kann, falls $\abs{\lambda}> 2 \cdot \norm{a}$ ist.
	Sei nun $\lambda_0 \in \mathbb{C}\setminus \sigma(a)$ und $\abs{\lambda-\lambda_0}< \norm{(a- \lambda_0)^{-1}}^{-1}$. Dann gilt 
	\[
		\norm{(\lambda-\lambda_0)(a - \lambda_0)^{-1}} \le \abs{\lambda-\lambda_0} \cdot \norm{(a-\lambda_0)^{-1}} <1   
	\]
	Also ist wieder nach \autoref{satz15_neumann} $\enbrace*{\ind- (\lambda-\lambda_0)(a-\lambda_0)^{-1}}^{-1} = \sum_{n=0}^{\infty} (\lambda-\lambda_0)^n (a-\lambda_0)^n$ und 
	\begin{align}
		-\sum_{n=0}^{\infty} (\lambda-\lambda_0)^n(a-\lambda_0)^{-(n+1)} &= (a- \lambda_0)^{-1} \enbrace*{(\lambda-\lambda_0)(a-\lambda_0)^{-1}-\ind}^{-1} \\
		&= \enbrace[\Big]{(\lambda-\lambda_0) \Underbrace{(a-\lambda_0)^{-1}(a-\lambda_0)}{=\ind}-(a- \lambda_0)}^{-1} = (\lambda-a)^{-1}
	\end{align}
	Für ein stetiges Funktional $\varphi \in A^*$ erhalten wir 
	\[
		\varphi\enbrace*{(\lambda-a)^{-1}} = \sum_{n=0}^{\infty} - \varphi \enbrace*{(a-\lambda_0)^{-(n+1)}} \cdot (\lambda-\lambda_0)^n
	\]
	falls $\abs{\lambda-\lambda_0} < \norm{(a- \lambda_0)^{-1}}^{-1}$. Daher ist $\lambda \mapsto \varphi \enbrace*{(\lambda-a)^{-1}}$ \Index{holomorph} auf 
	$\mathbb{C}\setminus \sigma(a)$. Falls $\sigma(a)= \emptyset$, so ist $\lambda \mapsto \varphi \enbrace*{(\lambda-a)^{-1}}$ \bet{ganz}\index{ganze Funktion}, d.h. holomorph 
	auf ganz $\mathbb{C}$. Wegen \eqref{19:eq:1} und dem Satz von Liouville\footnote{siehe \url{https://de.wikipedia.org/wiki/Satz_von_Liouville_\%28Funktionentheorie\%29}} ist 
	dann $\varphi \enbrace*{(\lambda-a)^{-1}}=0$ für alle $\lambda\in \mathbb{C}$. Aber $\varphi \in A^*$ war beliebig, und wegen Hahn-Banach gilt $(\lambda-a)^{-1}=0$. Dies 
	ist ein Widerspruch. 
\end{beweis}

\begin{satz}[{name=[Formel für den Spektralradius]},label=satz1.10]
	Sei $A$ eine unitale Banachalgebra. Für $a \in A$ gilt dann 
	\[
		r(a)=\lim_{n \to \infty}\norm{a^n}^{\frac{1}{n}}
	\]
\end{satz}
\begin{beweis}
	Sei $\varphi \in A^*$ ein stetiges Funktional. Wir haben im Beweis von \autoref{spec_nichtleer} gesehen, dass die Zuordnung 
	$\lambda \mapsto \varphi \enbrace*{(\lambda-a)^{-1}}$ analytisch ist für $\lambda \notin \sigma(a)$. Insbesondere gilt dies für $\abs{\lambda}> r(a)$. Dann ist auch
	$\lambda \mapsto \varphi \enbrace*{(\lambda^{-1}-a)^{-1}}$ analytisch für $0< \abs{\lambda}< \sfrac{1}{r(a)}$, ebenso 
	$\lambda \mapsto \varphi \enbrace*{(\ind- \lambda a)^{-1}}$. Die Singularität bei $0$ ist hebbar, also definiert $f(\lambda) := \varphi \enbrace*{(\ind-\lambda a)^{-1}}$
	eine analytische Funktion für $\abs{\lambda} < \sfrac{1}{r(a)}$. Dann existiert eine Folge $(\alpha_n)_{n \in N} \subset \mathbb{C}$ mit 
	$f(\lambda) = \sum_{n=0}^{\infty} \alpha_n \cdot \lambda^n$ für $\abs{\lambda}<\sfrac{1}{r(a)}$. Die Konvergenz ist absolut und gleichmäßig, also gilt insbesondere 
	$\abs{\alpha_n \cdot \lambda^n} \xrightarrow{\raisemath{-1pt}{n \to \infty}}0$. Falls andererseits $\abs{\lambda} < \sfrac{1}{\norm{a}} \le \sfrac{1}{r(a)}$ 
	ist, so gilt
	\[
		f(\lambda) = \varphi \enbrace*{(\ind- \lambda \cdot a)^{-1}} \StackText{\ref{satz15_neumann}}{=} \varphi \enbrace*{\sum_{n=0}^{\infty} (\lambda \cdot a)^n} = 
		\sum_{n=0}^{\infty} \varphi \enbrace[\big]{(\lambda \cdot a)^n} = \sum_{n=0}^{\infty} \varphi(a^n) \cdot \lambda^n
	\]
	also $\alpha_n =\varphi(a^n)$ für $n \in N$. Es gilt $\varphi(a^n) \cdot \lambda^n \xrightarrow{\raisemath{-1pt}{n \to\infty}} 0$. Aber $\varphi \in A^*$ war beliebig, 
	also ist für jedes $\varphi \in A^*$ die Folge $\enbrace*{\varphi(\lambda^n a^n)}_\mathbb{N} = \enbrace*{\varphi(a^n) \cdot \lambda^n}_{n \in \mathbb{N}}$ beschränkt für 
	$\abs{\lambda}<\sfrac{1}{r(a)}$.

	Nach dem Prinzip der gleichmäßigen Beschränktheit ist dann $(\lambda^n a^n)_{n \in \mathbb{N}}$ beschränkt für jedes $\lambda$ mit $\abs{\lambda}< \sfrac{1}{r(a)}$. Damit 
	existiert für $0 < \abs{\lambda}< \sfrac{1}{r(a)}$ ein $M \ge 0$ mit $\norm{\lambda^n a^n} \le M$, $n \in \mathbb{N}$, also 
	\[
		\norm{a^n}^{\sfrac{1}{n}}\le\frac{M^{\sfrac{1}{n}}}{\abs{\lambda}} \xrightarrow{n \to \infty} \frac{1}{\abs{\lambda}} 
	\]
	Es folgt $\limsup_n \norm{a^n}^{\sfrac{1}{n}} \le r(a)$. Falls $\lambda \in \sigma(a)$, so gilt nach \autoref{satz14_poly} $\lambda^n \in \sigma(a^n)$ und wegen 
	\autoref{spec_kompakt} $\abs{\lambda^n} \le \norm{a^n}$, also $\abs{\lambda} \le \norm{a^n}^{\sfrac{1}{n}}$ und es folgt insgesamt
	\begin{align}
		r(a) \le \inf_n \norm{a^n}^{\sfrac{1}{n}} \le \liminf_n \norm{a^n}^{\sfrac{1}{n}} \le \limsup_n \norm{a^n}^{\sfrac{1}{n}} \le r(a)
	\end{align}
	und damit $\lim_n \norm{a^n}^{\sfrac{1}{n}} = r(a)$, woraus die Behauptung folgt. 
\end{beweis}

\begin{proposition}[label=prop1.11_einfach,{name=[Einfache abelsche unitale Banachalgebra]}]
	Sei $A$ eine abelsche unitale Banachalgebra. Falls $A$ einfach\footnote{besitzt keine nichttrivialen Ideale} ist, so ist $A \cong \mathbb{C}$.\index{einfach} 
\end{proposition}
\begin{beweis}
	Sei $a \in A$, $\alpha \in \sigma(a)$. Dann ist $J := \overline{(a-\alpha \cdot \ind)\cdot A} \lhd A$ ein abgeschlossenes Ideal. Für jedes $b \in A$ ist $(a-\alpha)b$ nicht 
	invertierbar. Damit folgt dann $\norm{(a-\alpha)b - \ind}\ge 1$ für jedes $b \in A$ nach \autoref{satz15_neumann}. Damit ist auch 
	$\dist\enbrace*{(a-\alpha \cdot \ind)\cdot A,\ind}\ge 1$, also $\ind \notin J$ und damit $J \not= A$. Da $A$ einfach ist, folgt $J=\set{0}$, also muss 
	$a-\alpha \cdot \ind=0$ und weiter $a \in \mathbb{C} \cdot \ind$ gelten. Damit ist $A = \mathbb{C} \cdot \ind \cong \mathbb{C}$.
\end{beweis}

\begin{definition}[{name=[Charakter]}]
	Sei $A$ eine abelsche unitale Banachalgebra. Eine lineare, multiplikative, nicht verschwindende Abbildung $\chi \colon A \to \mathbb{C}$ heißt \Index{Charakter} auf $A$.  
	Es gilt dann $\chi(\ind)=1$. Wir setzen $\Omega(A) := \set[\big]{\chi \given \chi \text{ Charakter auf }A}$ (\Index{Charakterspektrum} oder einfach nur \Index{Spektrum} von 
	$A$). \todo{Revision 1}
\end{definition}

\begin{satz}[label=satz:1.13,{name=[Stetigkeit von Charakteren, Zusammenhang mit maximalen Idealen]}]
	Sei $A$ eine abelsche unitale Banachalgebra. Dann gilt
	\begin{enumerate}[(i)]
		\item Jeder Charakter auf $A$ ist stetig mit Norm $1$.
		\item Die Abbildung $\chi \mapsto \ker \chi$ definiert eine Bijektion  
		\[
			\Omega(A) \longrightarrow J_{\max} := \set[\big]{J \given J \subsetneq A \text{ ist ein maximales Ideal} }
		\]
	\end{enumerate}
\end{satz}
\begin{beweis}\leavevmode
	\begin{enumerate}[(i)]
		\item Sei $a \in A$ mit $\norm{a}<1$ und $\chi(a)=1$. Setze $b= \sum_{n=1}^{\infty} a^n$, dann ist $1 + b = \sum_{n=0}^{\infty} a^n = (1-a)^{-1}$. Also ist
		$1 = (1+b)(1-a) = 1 + b -a -ab$ und damit $ab+a=b$. Es folgt $\chi(a)\chi(b) + \chi(a)= \chi(b)$. Dies ist ein Widerspruch. Also muss $\norm{\chi}\le 1$ sein. Mit 
		$\chi(1)=1$ folgt $\norm{\chi}=1$.
		\item Für jedes $\chi \in \Omega(A)$ ist $\ker \chi \lhd A$ ein Ideal. Weiter ist der Kern abgeschlossen in $A$, da $\chi$ stetig ist. Wir zeigen, dass $\ker \chi$ maximal 
		ist. $\ker \chi \not= A$, denn $\chi \not= 0$. Falls $c \in A \setminus \ker \chi$, so gilt 
		\(
			\chi(c)^{-1} \cdot c-1 \in \ker \chi
		\).
		Also ist $\ker \chi + \mathbb{C} \cdot c = \ker \chi + \mathbb{C} \cdot 1=A$, da $a- \chi(a) \cdot 1 \in \ker \chi$ und damit 
		$a= \enbrace*{a- \chi(a) \cdot 1}+ \chi(a) \cdot 1 \in \ker \chi + \mathbb{C} \cdot 1$. Damit ist der Kern von $\chi$ maximal.
		
		Zur Injektivität: Es seien $\chi_1, \chi_2 \in \Omega(A)$ mit $\ker \chi_1 = \ker \chi_2$. Dann gilt für $a \in A$
		\begin{align}
			\chi_1(a) - \chi_2(a) = \chi_1 \enbrace[\Big]{\Underbrace{a - \chi_2(a) \cdot \ind}{\in \ker \chi_2=\ker \chi_1}} = 0
		\end{align}
		Also ist $\chi_1=\chi_2$. Für die Surjektivität betrachten wir ein maximales, nichttriviales Ideal $J \subsetneq A$. Dann ist $\dist \enbrace*{J,1}=1$, denn 
		\todo{Einsen auseinanderhalten} 
		$\mathring{B}_A(1,1) \subset \Inv(A)$ nach \autoref{satz15_neumann}. Dann ist $\ind \notin \overline{J} \lhd A$. Da $J$ maximal ist, folgt $J=\overline{J}$. Damit ist 
		$\sfrac{A}{J}$ eine einfache, abelsche, unitale Banachalgebra. Mit \autoref{prop1.11_einfach} folgt $\sfrac{A}{J} \cong \mathbb{C}$. Damit ist 
		$\chi \colon A \to \sfrac{A}{J} \cong \mathbb{C}$ ein Charakter mit $\ker \chi= J$. \qedhere
	\end{enumerate}
\end{beweis}

\begin{satz}[label=satz:1.14,{name=[Spektrum entspricht Menge der $\chi(a)$]}]
	Sei $A$ eine abelsche unitale Banachalgebra. Für $a \in A$ gilt $\sigma(a)=\set[\big]{\chi(a) \given \chi \in \Omega(A)}$.
\end{satz}
\begin{beweis}
	Für $\lambda \in \sigma(a)$ ist $J := (a -\lambda)A \lhd A$ ein echtes Ideal, da $(a-\lambda)$ nicht invertierbar ist. Nach dem Lemma von Zorn ist $J$ in einem maximalen echten 
	Ideal enthalten. Nach \autoref{satz:1.13} ist $J \subset \ker \chi$ für ein $\chi \in \Omega(A)$. Dann ist $\chi(a- \lambda)=0$ für einen Charakter $\chi$ und somit auch 
	$\chi(a)=\lambda$. Es gilt also $\sigma(a) \subseteq \set{\chi(a) \given \chi \in \Omega(A)}$.
	
	Für $\chi \in \Omega(A)$ gilt $a- \chi(a) \cdot \ind \in \ker \chi$. Damit kann $a- \chi(a) \cdot \ind$ nicht invertierbar sein. Also ist $\chi(a) \in \sigma(a)$ und es folgt 
	die zweite Inklusion.
\end{beweis}

\begin{satz}[label=satz:1.15,{name=[Charakterspektrum ist kompakt]}]
	Sei $A$ eine abelsche unitale Banachalgebra. Dann ist $\Omega(A)$ ein kompakter Hausdorffraum bezüglich der $\w^*$-Topologie.
\end{satz}
\begin{beweis}
	Es ist $\Omega(A) \subset B^1(A^*)$, wobei die Einheitskugel ein kompakt und Hausdorffsch ist (bezüglich der $\w^*$-Topologie) nach Banach-Alaoglu und da der Raum lokalkonvex 
	ist. Sei $\chi_\gamma \xrightarrow{\raisebox{-3pt}{$\w^*$}}\chi$ ein Netz. $\chi_\gamma$ ist multiplikativ für jedes $\gamma$ und $\chi_\gamma(\ind)=1$. Damit gelten diese Eigenschaften auch für 
	$\chi$.
\end{beweis}

\begin{definition}[{name=[Charakterspektrum und Gelfandtransformation]}]
	Sei $A$ eine abelsche unitale Banachalgebra. $\Omega(A)$ mit der $\w^*$-Topologie heißt \bet{Spektrum}\index{Spektrum!Charakterspektrum} von $A$. Für $a \in A$ definieren wir
	die \Index{Gelfandtransformation} $\hat{a} \colon \Omega(A) \to \mathbb{C}$ durch $\hat{a}(\chi)=\chi(a)$. 
\end{definition}

\begin{satz}[label=satz:117,{name=[Gelfandtransformation ist normvermindernd]}]
	Sei $A$ eine abelsche unitale Banachalgebra. Dann ist $\hat{\hspace{1ex}} \colon A \to C \enbrace*{\Omega(A)}$, $a \mapsto \hat{a}$ ein normvermindernder Homomorphismus und 
	$r(a) = \norm{\hat{a}}_{\infty}$ für $a \in A$. Weiter gilt $\sigma(a)= \hat{a}\enbrace*{\Omega(A)}$.
\end{satz}
\begin{beweis}
	$\hat{a} \colon \Omega(A) \to \mathbb{C}$ ist stetig nach Definition der $\w^*$-Topologie: Für $U \subset \mathbb{C}$ offen ist
	\[
		\hat{a}^{-1}(U) = \set[\big]{\chi \in \Omega(A) \given  \hat{a}(\chi) \in U} = \set[\big]{\chi \in \Omega(A) \given \chi(a) \in U} = 
		\set[\big]{\chi \in A^* \given \chi(a) \in U} \cap \Omega(A)
	\]
	Die linke Menge des Schnitts ist gleich
	\[
		\smashoperator[r]{\bigcup_{(\mu,\varepsilon), B(\mu,\varepsilon) \subset U}} \,\,\set[\big]{\varphi \in A^* \given \abs{\varphi(a)-\mu}<\varepsilon }
	\]
	und somit offen in $A^*$. Damit ist $\hat{a}^{-1}(U)$ offen in der Relativtopologie. Wir haben
	\begin{align}
		\sigma(a) \StackText{\ref{satz:1.14}}{=} \set[\big]{\chi(a) \given \chi \in \Omega(A)}  = \set[\big]{\hat{a}(\chi) \given \chi \in \Omega(A)} = \hat{a} 
		\enbrace[\big]{\Omega(A)}
	\end{align}
	Also 
	\[
		\norm{a} \StackText{\ref{spec_kompakt}}{\le} r(a) = \sup_{\lambda \in \sigma(a)} \abs*{\lambda} = \sup_{\chi \in \Omega(A)} \abs{\hat{a}(\chi)} 
		= \norm{\hat{a}}_{\infty} \qedhere
	\]
\end{beweis}
% section 1 (end)
\newpage

\section{$C^*$-Algebren und Funktionalkalkül} % (fold)
\label{sec:2}

\begin{definition}[{name=[Involutive Banachalgebra]}]
	Eine Banachalgebra heißt \Index{involutiv}, falls eine antilineare Abbildung ${}^* \colon A \to A$ existiert, für die folgenden Identitäten gelten:
	\[
		a^{**}=a \qquad \quad (ab)^*= b^* a^* \qquad \quad \norm{a^*} = \norm{a}  
	\] 
\end{definition}

\begin{definition}[{name=[$C^*$-Algebra]}]
	Eine \bet{$C^*$-Algebra}\index{C-Algebra@$C^*$-Algebra} ist eine involutive Banachalgebra mit $\norm{a^*a}= \norm{a}^2$ für $a \in A$.
\end{definition}

\begin{definition}[{name=[{selbstadjungiert, normal, unitär, Projektion, Isometrie}]}]
	Sei $A$ ein $C^*$-Algebra. $a \in A$ heißt 
	\[
		\begin{cases}
			\text{ \Index{selbstadjungiert} }  &\text{ falls } a=a^*\\
			\text{ \Index{normal} }  &\text{ falls }  aa^*=a^*a \\
			\text{ \Index{unitär} }  &\text{ falls }  aa^*=a^*a=\ind \\
			\text{ \Index{Projektion} }  &\text{ falls }   a=a^*=a^2 \\
			\text{ \Index{Isometrie} }  &\text{ falls } a^*a= \ind  
		\end{cases}
	\]
	Wo die Existenz eines Einselements nötig ist, fordern wir zusätzlich, dass $A$ unital ist.
\end{definition}

\begin{definition}[{name=[${ }^*$-Homomorphismus]}]
	Es seien $A,B$ zwei $C^*$-Algebren. Eine Abbildung $\pi \colon A \to B$ heißt \bet{${}^*$"=Homomorphismus}\index{Stern-Homomorphismus@Sten-, ${}^*$-Homomorphismus}, falls $\pi$ 
	linear, multiplikativ und ${}^*$-erhaltend ist. $\pi$ heißt
	\bet{${}^*$-Einbettung}\index{Stern-Einbettung@Stern-, ${}^*$-Einbettung}, falls $\pi$ zusätzlich isometrisch ist. 
\end{definition}

\begin{bemerkung}[{name=[Direkte Summe von $A$ mit den komplexen Zahlen]}]
	Falls $A$ eine (involutive) Banachalgebra ist, so ist $A \oplus \mathbb{C}$ mit offensichtlichen Vektorraumoperationen, der naheliegenden Multiplikation
	$(a,\lambda) \cdot (b,\mu) = \enbrace*{ab + \mu a + \lambda b, \lambda \mu}$ sowie $(a,\lambda)^* = (a^*,\lambda)$ im involutiven Fall und Norm definiert wir durch
	$\norm{(a,\lambda)}:= \norm{a}_{A} + \abs{\lambda}$ eine unitale (involutive) Banachalgebra mit $\ind=(0,1)$. $A \oplus \mathbb{C}$ ist jedoch keine $C^*$-Algebra, auch wenn
	$A$ eine $C^*$-Algebra ist.
\end{bemerkung}

\begin{satz}[label=satz:26,{name=[Unitalisierung]}]
	Sei $A$ eine $C^*$-Algebra. Dann existiert eine unitale $C^*$-Algebra $A^\simm$ mit einer isometrischen ${}^*$-Einbettung $\iota \colon A \hookrightarrow A^\simm$, sodass
	$\iota(A) \lhd A^\simm$ ein Ideal ist mit 
	\[
		\Faktor{A^\simm}{\iota(A)} \cong \begin{cases}
			0, &\text{ falls }A \text{ unital} \\
			\mathbb{C} , &\text{ sonst }  
		\end{cases}
	\]
\end{satz}
\begin{beweis}
	Für $a \in A$ definieren wir $\iota(a) \in \mathcal{B}(A)$, wobei $\mathcal{B}(A)$ die beschränkten Operatoren auf $A$ aufgefasst als Banachraum sind, durch
	$\iota(a)(b)=ab$. Definiere $A^\simm := \iota(A) + \mathbb{C} \cdot \ind_{\mathcal{B}(A)} \subset \mathcal{B}(A)$. Dann ist $A^\simm$ eine unitale Algebra mit Involution
	\[
		\enbrace*{\iota(a)+\lambda \cdot \ind_{\mathcal{B}(A)}}^* := \iota(a^*) + \overline{\lambda} \cdot \ind_{\mathcal{B}(A)}
	\]
	Dass $\iota(A) \lhd A^\simm$ ein Ideal ist, ist klar. $A^\simm$ ist normiert mittels der Operatornorm. Für $0 \neq a \in A$ gilt  
	\begin{align}
		\norm{a}_A = \frac{1}{\norm{a}_A}\norm{aa^*}_A = \norm*{a \enbrace*{\frac{1}{\norm{a^*} } \cdot a^*}}_A &\le \norm*{\iota(a) \enbrace*{\frac{1}{\norm{a^*}}\cdot a^* }}_A \\
		&\le \sup_{\norm{b}_A \le 1} \norm[\big]{\iota(a)(b)}_A = \norm[\big]{\iota(a)}_{\mathcal{B}(A)} \\
		&= \sup_{\norm{b}_A \le 1} \norm{ab}_A \\
		&\le \norm{a}_A   
	\end{align}
	Also ist $\iota$ isometrisch. Es folgt, dass $\iota(A) \subset \mathcal{B}(A)$ abgeschlossen ist, da $A$ vollständig ist. Weiter folgt, dass auch 
	$A^\simm = \iota(A) + \mathbb{C} \cdot \ind_{\mathcal{B}(A)}$ vollständig ist bezüglich $\norm{\cdot}_{\mathcal{B}(A)}$. Wir zeigen nun, dass $\norm{\cdot}_{\mathcal{B}(A)}$ 
	eine $C^*$-Norm ist auf $A^\simm$: 
	\begin{align}
		\norm*{\iota(a)+ \lambda \cdot \ind_{\mathcal{B}(A)}}_{\mathcal{B}(A)}^2 &= \sup_{b \in A, \norm{b} \le 1 } \norm*{ab + \lambda b}^2_A \\
		&= \sup_{b \in A, \norm{b} \le 1} \norm[\big]{\enbrace*{ab+\lambda b}^* (ab +\lambda b)}_A \\
		&= \sup_{b \in A, \norm{b} \le 1} \norm[\big]{b^* a^* ab + \lambda b^* a^* b + \overline{\lambda} b^* ab + \abs*{\lambda}^2 b^* b}_{A} \\
		&\le \sup_{b \in A, \norm{b} \le 1} \norm[\big]{a^*ab + \lambda a^* b + \overline{\lambda} a b + \abs{\lambda}^2 b }_A \tag{$*$} \label{eq:26:1}\\
		&= \norm[\big]{\iota \enbrace*{a^* a + \lambda a^* + \overline{\lambda} a} + \abs{\lambda}^2 \cdot \ind_{\mathcal{B}(A)} }_{\mathcal{B}(A)} \\
		&= \norm*{ \enbrace*{\iota(a) + \lambda \cdot \ind_{\mathcal{B}(A)}}^* \enbrace*{\iota(a) + \lambda \cdot \ind_{\mathcal{B}(A)}}}_{\mathcal{B}(A)} \\
		&\le \norm*{\enbrace*{\iota(a) + \lambda \cdot \ind_{\mathcal{B}(A)}}^*}_{\mathcal{B}(A)} \norm*{\iota(a) + \lambda \cdot \ind_{\mathcal{B}(A)}}_{\mathcal{B}(A)}   
		\tag{$**$} \label{eq:26:2}
	\end{align}
	Es folgt, dass 
	\begin{align}
		\norm*{\iota(a) + \lambda \cdot \ind_{\mathcal{B}(A)}}_{\mathcal{B}(A)} \le \norm*{\enbrace*{\iota(a) + \lambda \cdot \ind_{\mathcal{B}(A)}}^*}_{\mathcal{B}(A)}
		\le \norm*{\iota(a) + \lambda \cdot \ind_{\mathcal{B}(A)}}_{\mathcal{B}(A)} 
	\end{align}
	Damit folgt, dass \eqref{eq:26:1} und \eqref{eq:26:2} Gleichungen sind und $\norm*{\cdot }_{\mathcal{B}(A)}$ eine $C^*$-Norm auf $A^\simm$ ist. Falls $A$ unital ist, gilt
	$\iota(\ind_A) = \ind_{\mathcal{B}(A)}$ und somit $\iota(A)=A^\simm$. Ist $A$ nicht unital, so ist $A^\simm \cong \iota(A) \oplus \mathbb{C}$ als Vektorräume und es muss
	$\sfrac{A^\simm}{\iota(A)}\cong \mathbb{C}$ gelten. \todo{Revision 2}
\end{beweis}

\begin{bemerkung}[{name=[Alternative Unitalisierung]}]
	Man definiert oft auch eine unitale $C^*$-Algebra 
	\[
		A^+ := \begin{cases}
			A^\simm, &\text{ falls }A \text{ nicht unital} \\
			A  \oplus \mathbb{C} , &\text{ falls } A \text{ unital} 
		\end{cases}
	\]
	mit $\norm*{(a,\lambda)} := \max \set*{\norm*{a}_A, \abs{\lambda}}$. Dann ist in jedem Fall $A^+ \cong A \oplus \mathbb{C}$ als Vektorraum. $A^+$ lässt sich charakterisieren 
	durch die universelle Eigenschaft: 
	\[
		\begin{tikzcd}
			\mbox{ } & B \\
			A \rar[hook] \urar& A^+ \uar[dashed,"\exists !"']
		\end{tikzcd}
	\]
	wobei $B$ eine unitale $C^*$-Algebra ist und $A \to B$ ein ${}^*$-Homomorphismus ist. Eine ähnliche Beschreibung für $A^\simm$ ist möglich, aber etwas \enquote{subtiler}.
\end{bemerkung}

\begin{definition}[{name=[Spektrum in nichtunitalen $C^*$-Algebren]}]
	Sei $A$ eine nichtunitale $C^*$-Algebra. Für $a \in A$ setzen wir 
	\[
		\sigma(a) := \sigma_{A^\simm} \enbrace[\big]{\iota(a)} \quad \text{ und } \quad r(a) := r_{A^\simm} \enbrace[\big]{\iota(a)} 
	\]
	Wieder schreiben wir $\Omega(A) := \set[\big]{\chi \colon A \to \mathbb{C} \text{ Charaktere}}$ und \marginnote{Für $A$ unital sind maximale Ideale automatisch abgeschlossen}
	\[
		J_{\max}(A)= \set[\big]{J \given J \lhd A \text{ echtes maximales abgeschlossenes Ideal}}
	\]
\end{definition}

\begin{satz}[label=satz:29]
	Sei $A$ eine nichtunitale abelsche $C^*$-Algebra. Dann ist $\Omega(A)$ ein lokalkompakter Hausdorffraum bezüglich der $\w^*$-Topologie. Für $a \in A$ gilt 
	$\sigma(a) = \set*{\chi(a) \given \chi \in \Omega(A)} \cup \set*{0}$.
\end{satz}
\begin{beweis}
	Für $\chi \in \Omega(A)$ sei $\chi^\simm \colon A^\simm \to \mathbb{C}$ die eindeutig bestimmte, lineare, unitale Fortsetzung. Es gilt $\chi^\simm \in \Omega(A^\simm)$ und
	\[
		\Omega(A^\simm) = \set[\big]{\chi^\simm \given \chi \in \Omega(A)} \cup \set[\big]{\chi_\infty}
	\]
	wo $\chi_\infty \colon A^\simm \to \sfrac{A^\simm}{\iota(A)} \cong \mathbb{C}$. $\Omega(A^\simm)$ ist ein kompakter Hausdorffraum nach \autoref{satz:1.15}, also muss der linke 
	Teil der Vereinigung ein lokalkompakter Hausdorffraum sein.  Man überprüft, dass $\chi \mapsto \chi^\simm$ ein Homöomorphismus 
	$\Omega(A)\to \Omega(A^\simm) \setminus \set*{\chi_\infty}$ bezüglich der $\w^*$-Topologie ist. Es gilt
	\[
		\sigma(a) := \sigma_{A^\simm} \enbrace[\big]{\iota(a)} = \set[\big]{\chi \enbrace*{\iota(a)} \given \chi \in \Omega(A^\simm)} = 
		\Underbrace{\set[\big]{\chi \enbrace*{\iota(a)}\given\chi \in \Omega(A)}}{= \set*{\chi(a) \given \chi \in \Omega(A)}} \cup \Underbrace{\set[\big]{\chi_\infty 
		\enbrace*{\iota(a)}}}{=\set*{0}} \qedhere
	\]
\end{beweis}

\begin{proposition}[label=prop:210]
	Sei $A$ eine $C^*$-Algebra. Dann gilt
	\begin{enumerate}[(i)]
		\item Für eine Projektion $0 \neq p  \in A$ gilt $\norm*{p} =1$.
		\item Ist $\ind \in A$ und $u \in A$ unitär, dann ist $\norm*{u}=1$ und $\sigma(u) \subset \mathbb{T}$.
		\item Für $a \in A_{\sa}$ ist $r(a)= \norm*{a}$ und $\sigma(a) \subset \mathbb{R}$.
		\item Für $a \in A$ beliebig gilt $\norm*{a} = \norm*{a^* a}^{\sfrac{1}{2}}= r(a^* a)^{\sfrac{1}{2}}$.
	\end{enumerate}
\end{proposition}
\begin{beweis}
	\leavevmode
	\begin{enumerate}[(i)]
		\item Es gilt $\norm*{p} = \norm*{p^*p}= \norm*{p}^2$ und somit muss $\norm*{p}=1$ gelten.
		\item Es gilt $\norm*{u}^2= \norm*{u^* u}= \norm*{\ind}=1$, also $\norm*{u}=1=\norm*{u^*}$. Wenn $\lambda \in \sigma(u)$, dann gilt $\lambda \neq 0$ und es gilt
		$\lambda^{-1} \in \sigma(u^*)$. Da die Spektralradien von $u$ und $u^*$ kleiner gleich $1$ sind, muss $\abs{\lambda},\abs{\lambda^{-1}}\le 1$ gelten, also $\abs{\lambda}=1$
		und $\lambda \in \mathbb{T}$.
		\item Es gilt $\norm*{a^2}= \norm*{a^* a}= \norm*{a}^2$. Eine Induktion liefert $\norm*{a^{2^n}}= \norm*{a}^{2^n}$, also nach \autoref{satz1.10} 
		\[
			r(a) = \lim_{n \to \infty} \norm*{a^n}^{\sfrac{1}{n}} = \lim_{n \to \infty} \norm*{a^{2^n}}^{\sfrac{1}{2^n}} = \norm*{a}
		\]
		Weiter existiert\footnote{$(ia)^0 = \ind_{A^\simm}$ falls $A$ nicht unital} $e^{ia} := \sum_{n=0}^{\infty} \frac{1}{n!}(ia)^n \in A^\simm$ und ist unitär mit 
		\(
			\enbrace*{e^{ia}}^* = e^{(ia)^*} = e^{-ia} = \enbrace*{e^{ia}}^{-1}
		\), 
		also ist $\sigma(e^{ia}) \subset \mathbb{T}$ nach (ii). Falls $\lambda \in \sigma(a)$, dann setze $b := \sum_{n=1}^{\infty} \frac{1}{n!} i^n (a-\lambda)^{n-1}$. Dann gilt
		\[
			(a-\lambda) b e^{i \lambda} = \enbrace[\big]{e^{i(a-\lambda)}- \ind}e^{i \lambda} = e^{i a} - e^{i \lambda}
		\]
		Also ist $e^{i \lambda} \in \sigma(e^{ia}) \subset \mathbb{T}$ und somit $\lambda \in \mathbb{R}$.
		\item Folgt aus (iii). \qedhere
	\end{enumerate}
\end{beweis}

\begin{korollar}[label=kor:211]
	Es seien $A$ und $B$ zwei $C^*$-Algebren und $\pi \colon A \to B$ ein ${}^*$-Homomorphismus. Dann ist $\pi$ kontraktiv.
	Die Aussage bleibt richtig, falls $A$ nur eine involutive Banachalgebra ist. 
\end{korollar}
\begin{beweis}
	Wir dürfen $\pi$ als unital annehmen (betrachte sonst $\pi^{\hspace{-0.5pt}+}\colon A^{\hspace{-0.8pt}+} \to B^{\hspace{-0.5pt}\sim}$). Für $d \in A$ gilt 
	$\sigma_B \enbrace[\big]{\pi(d)} \subset \sigma_A(d)$, alos $r_B \enbrace[\big]{\pi(d)} \le r_A(d)$. Für $a \in A$ gilt
	\[
		\norm*{\pi(a)}^2 = \norm*{\pi(a)^* \pi(a)} = \norm*{\pi(a^*a)} = r_B \enbrace[\big]{\pi(a^*a)} \le r_A(a^*a) = \norm*{a}^2 \qedhere
	\]
\end{beweis}

\begin{korollar}
	Auf einer ${}^*$-Algebra gibt es höchstens eine $C^*$-Norm.
\end{korollar}
\begin{beweis}
	Sei $A$ eine ${}^*$-Algebra mit $C^*$-Normen $\norm*{\cdot}_1$ und $\norm*{\cdot}_2$. Wir bezeichnen die $C^*$-Algebren entsprechend mit $A_1$ und $A_2$. Sei 
	$\pi_1 \colon A_1 \to A_2$ und $\pi_2 \colon A_2 \to A_1$ die Identitäten. Dann gilt
	\[
		\norm*{a}_1 = \norm*{\pi_2(a)}_1 \le \norm*{a}_2 = \norm*{\pi_1(a)}_2 \le \norm*{a}_1 \qedhere
	\]
\end{beweis}

\begin{korollar}[label=kor:213]
	Sei $A$ eine $C^*$-Algebra und $\chi$ ein Charakter auf $A$. Dann gilt $\chi(a^*) = \overline{\chi(a)}$ für $a \in A$, das heißt $\chi$ ist ein ${}^*$-Homomorphismus.
\end{korollar}
\begin{beweis}
	Sei $a \in A$, dann ist $a=\frac{1}{2}(a+a^*) + i \cdot \frac{1}{2i}(a-a^*) =: b + i \cdot c$ und $b,c$ sind selbstadjungiert. Es gilt 
	\[
		\chi(b) \StackText{\ref{satz:29}}{\in} \sigma(b) \StackText{\ref{prop:210}(iii)}{\subset} \mathbb{R} \quad \text{ und } \quad \chi(c) \in \sigma(c) \subset \mathbb{R} 
	\]
	Damit folgt $\chi(a^*) = \chi(b - i \cdot c) = \chi(b) - i \cdot \chi(c) = \overline{\chi(b) + i \cdot \chi(c)} = \overline{\chi(a)}$.
\end{beweis}

Wir erhalten eine Version von \autoref{satz:1.13} für nicht notwendigerweise unitale $C^*$-Algebren:

\begin{korollar}
	Sei $A$ eine abelsche $C^*$-Algebra. Dann gilt
	\begin{enumerate}[(i)]
		\item Jeder Charakter $\chi \colon A \to \mathbb{C}$ ist stetig mit Norm $\le 1$.
		\item Die Abbildung $\chi \mapsto \ker \chi$ definiert eine Bijektion 
		\[
			\Omega(A) \to J_{\max} = \set[\big]{J \given J \subsetneq A \text{ maximales, abgeschlossenes Ideal}}
		\]
	\end{enumerate}
\end{korollar}
\begin{beweis}
	\leavevmode
	\begin{enumerate}[(i)]
		\item folgt aus \autoref{kor:211} und \autoref{kor:213} oder auch aus \autoref{satz:1.13} und \autoref{satz:26}.
		\item Wir zeigen nur die Surjektivität: Sie $J \subsetneq A$ ein abgeschlossenes, maximales Ideal. Dann ist $J \subsetneq A^\simm$ ein abgeschlossenes Ideal. 
		\begin{align}
			J \subset A \cap \bigcap_{\substack{I \lhd A^\simm \text{ max} \\ J \subset I} } \hspace{-3pt} I  
			\StackText{\ref{satz:1.13}}{=} A \cap \bigcap_{\substack{\chi \in \Omega(A^\simm) \\ J \subset \ker \chi}} \ker \chi 
			&\StackText{\ref{satz:29}}{=} A \cap \ker \chi_\infty \cap \bigcap_{\substack{\chi \in \Omega(A) \\ J \subset \ker \chi^\simm}} \ker \chi^\simm  
		\end{align}
		Damit ist $J$ enthalten im Schnitt über alle Kerne der $\chi \in \Omega(A)$, wobei $J \subset \ker \chi$. Damit ist $J \subset \ker \chi$ für ein $\chi \in \Omega(A)$ und
		da $\ker \chi \lhd A$ abgeschlossen und maximal ist, folgt $J=\ker \chi$. \qedhere
	\end{enumerate}
\end{beweis}

\begin{beispiel}
	\begin{enumerate}[(i)]
		\item Sei $\mathcal{H}$ ein Hilbertraum. Dann ist $\mathcal{B}(\mathcal{H})$ eine $C^*$-Algebra mit Operatornorm und Adjunktion: Für $T \in \mathcal{B}(\mathcal{H})$ gilt
		\[
			\norm*{T} = \sup \set[\big]{\abs*{\skal*{\eta}{T \xi}} \given \norm*{\xi}, \norm*{\eta} \le 1}
		\]
		also
		\begin{align}
			\norm*{T^* T} = \sup \set[\big]{\abs*{\skal*{\eta}{T^* T \xi}} \given \norm*{\xi}, \norm*{\eta} \le 1} &\ge \sup \set[\big]{ \abs*{\skal*{T \eta}{T \eta}} \given 
			\norm*{\eta} \le 1}\\
			&= \enbrace*{\sup \set[\big]{\norm*{T \eta} \given \norm*{\eta} \le 1}}^2 = \norm*{T}^2 \ge \norm*{T^*T} \qedhere
		\end{align}
		\item Insbesondere ist $M_n := M_n(\mathbb{C}) = \mathcal{B}(\mathbb{C}^n)$ eine $C^*$-Algebra für $n \in \mathbb{N}$.
		\item Jede abgeschlossene ${}^*$-Unteralgebra von $\mathcal{B}(\mathcal{H})$ ist eine $C^*$-Algebra, zum Beispiel $\mathcal{K}(\mathcal{H})$.
		\item Sei $X$ ein lokalkompakter Hausdorffraum. Dann ist $C_0(X)$ eine $C^*$-Algebra mit punktweisen Operationen, Involution $f^*= \overline{f}$ und der Supremumsnorm.
		Tatsächlich ist jede abelsche $C^*$-Algebra von dieser Form: 
	\end{enumerate}
\end{beispiel}

\begin{satz}[Gelfand,label=satz16:gelfand]
	Sei $A \neq 0$ eine abelsche $C^*$-Algebra. Dann ist $\hat{\hspace{1ex}} \colon A \to C_0 \enbrace*{\Omega(A)}$\marginnote{$C_0 \enbrace*{\Omega(A)} = C \enbrace*{\Omega(A)}$, 
	falls $A$ unital} ein isometrischer ${}^*$-Isomorphismus.
\end{satz}
\begin{beweis}
	Betrachte $a \mapsto  \enbrace*{\chi \mapsto \chi(a), \chi \in \Omega(A) \cup \set*{\chi_\infty}}$ in folgendem Diagramm
	\[
		\begin{tikzcd}
			A \dar[hook,"\iota"] \rar["\hat{\hspace{1ex}}"] 
			& \set[\big]{f \in C \enbrace*{\Omega(A) \cup\set*{\chi_\infty}}\given f(\chi_\infty)=0} = C_0\enbrace*{\Omega(A)} \dar[hook] \\
			A^\simm \rar["\hat{\hspace{1ex}}"] & C \enbrace[\big]{\Omega(A^\simm)} = C \enbrace[\big]{ \set[\big]{\chi^\simm \given \chi \in \Omega(A)} \cup \set*{\chi_\infty}} 
			\cong C \enbrace[\big]{\Omega(A) \cup \set*{\chi_\infty}} 
		\end{tikzcd}
	\]
	Die Gelfandtransformation ist ein ${}^*$-Homomorphismus nach \autoref{kor:213} und \autoref{satz:117}. Weiter gilt 
	\[
		\norm*{a}^2 = \norm*{a^*a} \StackText{\ref{prop:210}}{=} r(a^*a) \StackText{\ref{satz:117}}{=} \norm[\big]{\widehat{a^*a}}_\infty = \norm*{\hat{a^*} \hat{a}}_\infty
		= \norm*{\hat{a} \hat{a}}_\infty = \norm[\big]{\abs{\hat{a}}^2 }_\infty = \norm*{\hat{a}}^2_\infty
	\]
	Damit ist $a \mapsto \hat{a}$ isometrisch. Damit folgt, dass $C := \set*{\hat{a} \given a \in A} \subset C_0 \enbrace*{\Omega(A)}$ abgeschlossen ist. Außerdem ist $C$ eine
	${}^*$-Unteralgebra, also insbesondere invariant unter komplexer Konjugation. Weiter trennt $C$ die Punkte von $\Omega(A)$ und zu jedem $\chi \in \Omega(A)$ existiert ein 
	$a \in A$ mit $\hat{a}(\chi)= \chi(a) \neq 0$, das heißt $\ev(C) \neq \set*{0}$ für alle $\chi \in \Omega(A)$. Mit Stone-Weierstraß folgt, dass 
	$C \subset C_0 \enbrace*{\Omega(A)}$ dicht ist. Da $C$ abgeschlossen ist, folgt also $C=C_0 \enbrace*{\Omega(A)}$.
\end{beweis}

\begin{bemerkung}[label=bem217]
	Sei $X$ ein kompakter Hausdorffraum und $f \in C(X)$, dann ist $\sigma_{C(X)}(f)=f(X)$.
\end{bemerkung}
\begin{beweis}
	Sei $\delta \colon X \to \Omega\enbrace*{C(X)}$ definiert durch $x \mapsto \enbrace*{f \mapsto f(x)}$ ist ein Homöomorphismus wegen $C(X) \cong C \enbrace[\big]{\Omega(C(X))}$. 
	Es gilt
	\begin{align}
		\sigma_{C(X)}(f) \StackText{\ref{satz:117}}{=} \hat{f} \enbrace[\big]{\Omega(C(X))} = \set*{\hat{f}(\chi) \given \chi \in \Omega \enbrace[\big]{C(X)}} 
		&= \set*{\chi(f) \given \chi \in \Omega \enbrace[\big]{C(X)}} \\
		&= \set[\big]{\delta(x)(f) \given x \in X} = f(X) \qedhere
	\end{align}
\end{beweis}

\begin{definition}
	Sei $S \subset A$ Teilmenge einer $C^*$-Algebra $A$. Dann bezeichnet $C^*(S) \subset A$ die kleinste $C^*$-Unteralgebra von $A$, die $S$ enthält; also
	\[
		C^*(S) = \smashoperator[r]{\bigcap_{\substack{B \subset A \text{ Unteralg.}\\ S \subset B}}} \enspace B
	\]
\end{definition}

\begin{satz}[Funktionalkalkül]
	Sei $A$ eine $C^*$-Algebra und $a \in A$ normal. Dann existiert ein eindeutig bestimmter ${}^*$-Isomorphismus $C^*(a) \cong C_0 \enbrace[\big]{\sigma(a)\setminus \set*{0}}$ 
	durch $a \mapsto \id_{\sigma(a)\setminus \set*{0}}$.
	Falls $A$ unital ist, so existiert ein eindeutig bestimmter unitaler ${}^*$-Isomorphismus $C^*(a,\ind) \cong C \enbrace*{\sigma(a)}$ durch $a \mapsto \id_{\sigma(a)}$.
	
	Für $f \in C_0\enbrace*{\sigma_A(a) \setminus \set*{0}}$ (bzw. $f \in C \enbrace*{\sigma_A(a)}$) bezeichnen wir mit $f(a) \in A$ das Urbild von $f$ unter diesem Isomorphismus.
	Es gilt $f(a) \in C^*(a,\ind_{A^\simm}) \subset A^\simm$, falls $f(0)\neq 0$. Weiter gilt 
	\[
		\sigma_A \enbrace[\big]{f(a)} = f \enbrace[\big]{\sigma_A(a)}
	\]
	Für $g \in C \enbrace*{f(\sigma_A(a))}$ gilt dann $g \enbrace*{f(a)}= (g \circ f)(a)$.
\end{satz}
\begin{beweis}
	Sei zunächst $A$ unital. Da $a \in A$ normal ist, ist $C^*(a,\ind) = \overline{^* \Alg(a,\ind)} \subset A$ abelsch. Die Abbildung
	\mapdef{\alpha \colon\Omega \enbrace*{C^*(a,\ind)}}{\hat{a} \enbrace*{\Omega \enbrace*{C^*(a,\ind)}}\StackText{\ref{satz:117}}{=} \sigma_{C^*(A,\ind)}\subset \mathbb{C}}%
	{\chi}{\hat{a}(\chi)=\chi(a)}{}
	ist ein Homöomorphismus: Die Stetigkeit ist klar, denn $\Omega \enbrace*{C^*(a,\ind)}$ trägt die $\w^*$-Topologie. Die Surjektivität ist trivial. Injektivität: Falls 
	$\chi_1(a)= \hat{a}(\chi_1) = \hat{a}(\chi_2) = \chi_2(a)$ ist, so ist $\chi_1|_{^*\Alg(a,\ind)} = \chi_2|_{^*\Alg(a,\ind)}$ und mit der Stetigkeit von $\chi_1$ und $\chi_2$ 
	folgt, dass $\chi_1=\chi_2$ sein muss.\todo{Abstand * und Alg fixen} Die Inverse ist stetig, da $\Omega \enbrace*{C^*(a,\ind)}$ kompakt ist.
	
	Wir erhalten nach \autoref{satz16:gelfand} $C^*(a,\ind) \cong C \enbrace*{\Omega(C^*(a,\ind))}$ mittels $a \mapsto \hat{a}$. und 
	$C \enbrace*{\Omega(C^*(a,\ind))} \cong C \enbrace*{\sigma_{C^*(a,\ind)}(a)}$ mittels $\hat{a} \mapsto \hat{a} \circ \alpha^{-1}$. Für 
	$\alpha(\chi)=t \in \sigma_{C^*(a,\ind)}(a)$ haben wir 
	\[
		\enbrace*{\hat{a} \circ \alpha^{-1}}(t) = \enbrace*{\hat{a} \circ \alpha^{-1}} \enbrace*{\alpha(\chi)} = \hat{a}(\chi) = \alpha(\chi)=t
	\]
	Also ist $\hat{a} \circ \alpha^{-1}=\id_{\sigma_{C^*(a,\ind)}(a)}$. Wir nenen die Komposition der beiden Isomorphismen $\beta$.\todo{schöner machen}
	
	Für $f \in C \enbrace*{\sigma_{C^*(a,\ind)}(a)}$ gilt 
	\[
		\sigma_{C^*(a,\ind)} \enbrace*{f(a)} = \sigma_{C^*(a,\ind)} \enbrace*{\beta^{-1}(f)} = \sigma_{\beta \enbrace*{C^*(a,\ind)}}(f) = \sigma_{C(\sigma_{C^*(a,\ind)}(a))} (f)
		\StackText{\ref{bem217}}{=} f \enbrace*{\sigma_{C^*(a,\ind)}(a)}
	\]
	Falls $A$ nicht unital ist, so haben wir\todo{hübsch machen bzw. leserlich}
	\[
		\begin{tikzcd}
			C^*(a,\ind_{A^\simm}) \rar["\StackText{\ref{satz16:gelfand}}{\cong}"] & C \enbrace*{\Omega \enbrace*{C^*(a,\ind_{A^\simm})}} \rar["\StackText{s.o.}{\cong}"]
			& C \enbrace*{\sigma_{C^*(a,\ind)}(a)} \\
			C^*(a) \rar["\StackText{\ref{satz16:gelfand}}{\cong}"] 
			& C_0 \enbrace*{\Omega \enbrace*{C^*(a)}} \StackText{\ref{satz:29}}{=} C_0 \enbrace*{\Omega \enbrace*{C^*(a,\ind_{A^\simm})} \setminus \set*{\chi_\infty}} \rar[two heads]
			& C_0 \enbrace*{\sigma_{C^*(a,\ind_{A})(a) \setminus \set*{0}}} =\set*{f \in C \enbrace*{\sigma_{C^*(a,\ind_{A^\simm})}(a)} \given f(0)=0} \\
			a \rar[mapsto] & \hat{a} \rar[mapsto] & \id_{\sigma_{C^*(a,\ind_{A^\simm})}(a)}
		\end{tikzcd}
	\]
	Der Funktionalkalkül $\beta^{-1}$ ist ein ${}^*$-Homomorphismus, daher gilt für ${}^*$-Polynome $g$, dass 
	\[
		g \enbrace*{f(a)} = g \enbrace*{\beta^{-1}(f)}= \beta^{-1}(g \circ f) = (g \circ f) (a)
	\]
	Für $g\in C \enbrace*{\sigma_{C^*(a,\ind)}(a)}$ approximieren wir $g$ gleichmäßig auf $\sigma_{C^*(a,\ind)}(a)$ durch $^*$-Polynome und benutzen die Stetigkeit von $\beta^{-1}$.
	Es bleibt zu zeigen, dass $\sigma_{C^*(a,\ind)}(a)= \sigma_A(a)$ ist. Offensichtlich gilt $\Inv \enbrace*{C^*(a,\ind_A)} \subseteq  \Inv(A)$ und somit
	$\sigma_A(a) \subseteq \sigma_{C^*(a,\ind_A)}(a)$. Um zu zeigen, dass auch die andere Inklusion gilt, sei $b \in C^*(a,\ind_A)$ mit $\norm*{b}\le 1$ und $b \in \Inv(A)$.
	Wir zeigen, dass $b$ normal ist: Es sind $b,b^* \in C^*(a,\ind_A)$ und  $C^*(b,b^{-1})$ ist abelsch, da $(b^{-1})^* = (b^*)^{-1}$ und 
	\[
		b(b^*)^{-1} = b(b^*)^{-1} b^{-1} b = b \enbrace*{bb^*}^{-1} b = b(b^*b)^{-1}b = (b^*)^{-1}b
	\]
	$\norm*{b^*b} \le 1$ und $\ind - b^*b$ ist selbstadjungiert und $0<\chi(b^*b) \le 1$ für $\chi \in \Omega \enbrace*{C^*(b,b^{-1})}$. Damit ist
	\[
		\sigma_{C^*(b,b^{-1})} \enbrace*{\ind - b^*b} \StackText{\ref{satz:117}}{=} \set*{\chi \enbrace*{\ind - b^*b} \given \chi \in \Omega \enbrace*{C^*(b,b^{-1})}} \subseteq 
		[0,1)
	\]
	Da das Spektrum $\sigma_{C^*(b,b^{-1})}(b)$ kompakt ist, folgt $1 > r(\ind - b^*b)= \norm*{\ind -b^*b}$. Es foglt
	\[
		(b^*b)^{-1} = \sum_{n=0}^{\infty} (\ind-b^*b)^n \in C^*(b,\ind)
	\]
	Damit folgt $b^{-1}=(b^*b)^{-1}b^* \in C^*(b,\ind) \subseteq C^*(a,\ind)$.
	
	Wenn $b \in C^*(a,\ind)$ invertierbar in $A$ ist, so ist $b^{-1} \in C^*(a,\ind)$, womit die andere Inklusion $\sigma_{C^*(a,\ind)}(a) \subseteq \sigma_A(a)$ folgt.
\end{beweis}

\begin{korollar}[label=korr:220]
	Sei $A \subseteq B$ eine $C^*$-Algebra und $a \in A$. Dann gilt
	\[
		\sigma_A(a) \cup \set*{0} = \sigma_B(a) \cup \set*{0}
	\]
	\begin{enumerate}[1)]
		\item Wenn $A$ und $B$ unital sind und $\ind_A=\ind_B$ ist, dann gilt $\sigma_A(a) = \sigma_B(a)$.
		\item Wenn $A$ nicht unital ist und $B$ unital ist, so ist $0 \in \sigma_A(a) = \sigma_B(a)$. 
		\item Wenn beide unital sind, die Einselemente aber nicht übereinstimmen, so gilt $\sigma_A(a) \cup \set*{0} = \sigma_B(a)$.
	\end{enumerate}
\end{korollar}
\begin{beweis} \leavevmode
	\begin{enumerate}[1)]
		\item Wenn $b$ normal ist, gehen wir genau wie im vorigen Beweis vor. Andernfalls betrachten wir $b^*b \in \Inv(B)$. Dann gilt
		\[
			(b^*b)^{-1} b^* = b^{-1} \in A
		\]
		\item Betrachte $A^\simm \cong C^*(a,\ind_B) \subseteq B$. Dann gilt
		\[
			\sigma_A(a) := \sigma_{A^\simm}(a) \StackText{1)}{=}\sigma_B(a)
		\]
		\item Wenn $\ind_A \neq \ind_B$, so gilt für $a \in \Inv(B)$ auch $a^{-1}\in A$. Dann folgt aber $\ind_B = a^{-1}a\in A$. Dies ist ein Widerspruch, also muss 
		$0 \in \sigma_B(a)$ gelten. Wenn $\lambda \neq 0$, dann gilt $\lambda \ind_A - a \in \Inv(A)$ genau dann, wenn $\lambda \ind_B -a \in \Inv(B)$. Dies sieht man wie folgt:
		\begin{align}
			\enbrace*{\lambda \ind_B- a}^{-1} = \enbrace*{\lambda \ind_A -a}^{-1} + \lambda^{-1}\enbrace*{\ind_B-\ind_A}
		\end{align}
		Damit ist $\sigma_A(a) \setminus \set*{0} = \sigma_B(a) \setminus \set*{0}$. \qedhere
	\end{enumerate}
\end{beweis}
% section 2 (end)
\newpage

\section{Positivität} % (fold)
\label{sec:3}

\begin{definition}[{name=[Positivität]}]
	Sei $A$ eine $C^*$-Algebra. Ein Element $a \in A$ heißt \Index{positiv} ($a \ge 0$), wenn $a=a^*$ und $\sigma(a) \subseteq [0,\infty)$ ist. Mit $A_+$ bezeichnen wir die 
	Menge der positiven Elemente. Wenn $b,c \in A_{\sa}$ sind, dann definieren wir \todo{Revision 3}
	\[
		b \le c \enspace :\Longleftrightarrow \enspace c-b \ge 0
	\]
\end{definition}

\begin{bemerkung} \leavevmode
	\begin{enumerate}[(i),itemsep=0pt]
		\item Für $a \in A \subseteq B$ gilt $a \ge 0$ in $A$ genau dann, wenn $a \ge 0$ in $B$ gilt (nach \autoref{korr:220}).
		\item Für $a \in A = C_=(X)$ ist $a \ge 0$ genau dann, wenn $a(x) \ge 0$ für alle $x \in X$ gilt (nach \ref{bem217}).
	\end{enumerate}
\end{bemerkung}

\begin{proposition}[label=prop:33]
	Sei $A$ eine $C^*$-Algebra und $a \in A_+$. Dann gibt es ein eindeutiges $b \in  A_+$, sodass $b^2=a$. Wir schreiben $a^{\sfrac{1}{2}}$ für dieses Element $b$.
\end{proposition}
\begin{beweis}
	Ein solches $b \in C^*(a) \subseteq A$ existiert nach dem Funktionalkalkül, da $C^*(a) \cong C_0 \enbrace*{\sigma(a) \setminus \set*{0}}$ via $a \mapsto \id_{\sigma(a)} \ge 0$.
	Sei $c \in A_+$ ein weiteres Element mit $c^2=a$. Dann gilt $ca=c^3=ac$. Damit ist $C^*(a,c)$ abelsch. Sei $\Omega := \Omega \enbrace*{C^*(a,c)}$. Sei 
	$\Phi \colon C^*(a,c) \to C_0(\Omega)$ die Gelfandtransformation. Dann sind $\Phi(b)$ und $\Phi(c)$ Quadratwurzeln von $\Phi(a)$. Anders ausgedrückt gilt für alle 
	$\omega \in \Omega$
	\[
		\Phi(b)(\omega)^2 = \Phi(a)(\omega) = \Phi(c)(\omega)^2
	\]
	Damit folgt $\Phi(b)(\omega)=\Phi(c)(\omega)$ für alle $\omega$ und somit $\Phi(b)=\Phi(c)$, also $b=c$.
\end{beweis}

\begin{proposition}[label=prop:34]
	Sei $a$ eine unitale $C^*$-Algebra, $a \in A_{\sa}$, $t \in \mathbb{R}_+ = [0,\infty)$. Dann gilt
	\begin{enumerate}[(i),itemsep=0pt]
		\item Aus $\norm*{a-t}\le t$ folgt $a \ge 0$
		\item Aus $a \ge 0$ und $\norm*{a}\le t$ folgt $\norm*{a-t}\le t$
	\end{enumerate}
\end{proposition}
\begin{beweis}
	Wir können annehmen, dass $A=C^*(a,\ind)=C \enbrace*{\sigma(a)}$ gilt.
	\begin{enumerate}[(i)]
		\item Es gilt $\sup_{x \in \sigma(a)} \abs*{x-t} = \norm*{a-t} \le t$, woraus $\abs*{x-t} \le t $ für alle $x \in \sigma(a)\subseteq \mathbb{R}$ folgt. Damit folgt
		$x \ge 0$ für alle $x \in \sigma(a)$, also $a \ge 0$.
		\item Aus $\norm*{a}\le t$ und $a \ge 0$ folgt $0 \le x \le t$ für alle $x \in \sigma(a)$. Subtrahieren mit $t$ liefert $-t \le x-t \le 0$ für alle $x$, also 
		$\abs*{x-t}\le t $. Damit folgt wieder $\norm*{a-t}\le t$. \qedhere
	\end{enumerate}
\end{beweis}

\begin{korollar}[label=korr:35]
	Sei $A$ eine $C^*$-Algebra. Dann ist $A_+$ ein Kegel, das heißt für $a,b \in A_+$ und $\lambda, \mu \in \mathbb{R}_+$ gilt $\lambda \cdot a + \mu \cdot b \in A_+$
\end{korollar}
\begin{beweis}
	Wir dürfen annehmen, dass $A$ unital ist. Dass $\lambda a \ge 0$ und $\mu b  \ge 0$ gilt, ist klar. Wir müssen also nur zeigen, dass aus $a,b\in A_+$ $a+b\in A_+$ folgt. Es gilt
	\begin{align}
		\norm*{a- \norm*{a}\cdot \ind_A} \le \norm*{a} \quad \text{ und } \quad \norm*{b- \norm*{b} \cdot \ind_A} 
	\end{align}
	nach voriger \autoref{prop:34}(ii). Damit folgt weiter
	\[
		\norm*{a+b - \enbrace*{\norm*{a} + \norm*{b}}} \le \norm*{a - \norm*{a}} + \norm*{b - \norm*{b}} \le \norm*{a} + \norm*{b}
	\]
	Also ist $a+b \ge 0$ nach \autoref{prop:34}(i). 
\end{beweis}

\begin{satz}[label=satz:36]
	Sei $A$ eine $C^*$-Algebra und $a \in A$. Dann gilt $a^*a \ge 0$.
\end{satz}
\begin{beweis}
	Definiere Funktionen $f,g,h \colon \mathbb{R} \to \mathbb{R}$ durch
	\missingfigure{Graphen der Funktionen}
	Es gilt $a^*a \in A_{\sa}$, $a^*a = f(a^*a)- g(a^*a)$ und $f(a^*a) g(a^*a)=0$. Für $c=a g(a^*a)$ ergibt sich
	\begin{align}
		-c^*c = -g(a^*a) a^*a g(a^*a) = -g(a^*a) \enbrace[\big]{f(a^*a) - g(a^*a)} g(a^*a) = g(a^*a)^3 \in A_+
	\end{align}
	Nach \autoref{bem:13} gilt $\sigma(-cc^*) \setminus \set*{0} = \sigma(-c^*c) \setminus \set*{0}$, also $-cc^* \in A_+$. Schreibe nun $c=d + i e$ mit $d,e \in A_\sa$. Dann gilt
	\begin{align}
		c^*c+cc^* = 2d^2+ 2e^2
	\end{align}
	Also ist $c^*c= 2d^2+2e^2- cc^* \ge 0$. Damit ist $\sigma(c^*c) \subseteq \mathbb{R}_+ \cap \mathbb{R}_- = \set*{0}$. Also ist $\norm*{c}^2= \norm*{c^*c}= r(c^*c)=0$. Damit ist
	$g(a^*a)^3=0$ und $g(a^*a)=0$ and $a^*a=f(a^*a)\ge 0$.
\end{beweis}

\begin{satz}
	Sei $A$ eine $C^*$-Algebra. Dann gilt
	\begin{enumerate}[(i)]
		\item $A_+ = \set*{a^*a \given a \in A}$
		\item $A = \Span A_+$
		\item $A_+$ ist ein abgeschlossener Kegel in $A$
		\item $\le$ ist eine partielle Ordnung auf $A_\sa$
		\item Für $a,b \in A_\sa$ mit $a \le b$ gilt $cac* \le cbc^*$.
		\item Für $0 \le a \le b$ gilt $\norm*{a} \le \norm*{b}$
		\item Wenn $A$ unital ist und $a \in A_+ \cap \Inv(A)$, dann folgt aus $a \le b$ stets $b \in \Inv(A)$ und $0 \le b^{-1} \le a^{-1}$.
	\end{enumerate}
\end{satz}
\begin{beweis}
	\leavevmode
	\begin{enumerate}[(i)]
		\item Folgt aus \autoref{satz:36} und \autoref{prop:33}.
		\item Wenn $a \in A$ ist, dann können wir $a=b+ic$ mit $b,c \in A_\sa$ schreiben. Definiere $b_+ := f(b)$, $b_- := g(b)$, $c_+ := f(c)$, $c_- := g(c)$. Diese liegen alle in 
		$A_+$. Dann gilt $b=b_+ - b_-$, $c=c_+ -c_-$ und $a= b_+ -b_- + i (c_+ -c_-) \in \Span A_+$.
		\item Kegel ist klar. Zur Abgeschlossenheit: Sei $a \in \overline{A}^{\norm*{\cdot }}$. Dann gilt $a=a*$, da $A_\sa$ abgeschlossen ist. Für jedes $\varepsilon>0$ gibt es
		$b \in A_+$ mit $\norm*{b-a} < \varepsilon$. Dann gilt
		\begin{align}
			\norm*{(a+\varepsilon-b) - \varepsilon} \le \varepsilon
		\end{align}
		$a+ \varepsilon-b \ge 0$ nach \autoref{prop:34}. Dann ist $a+\varepsilon= a+ \varepsilon -b +b \ge 0$ nach \autoref{korr:35}. Es folgt 
		$\sigma(a+\varepsilon) \subseteq [0,\infty)$ und damit $\sigma(a) \subseteq [\varepsilon,\infty)$. Lassen wir $\varepsilon$ gegen $0$ gehen, so folgt 
		$\sigma(a) \subseteq [0,\infty)$. Es folgt $a \ge 0$ und damit $\overline{A_+} =A_+$.
	\end{enumerate}
\end{beweis}
% section 3 (end)





\cleardoubleoddemptypage
\pagenumbering{Alph}
\setcounter{page}{1}

\printindex
\listoffigures
\todototoc
\listoftodos[To-do's und andere Baustellen]
\end{document}